\documentclass[11pt, a4paper]{article}

\usepackage[utf8]{inputenc}
\usepackage[T1]{fontenc}
\usepackage{amsmath, amssymb, amsthm, mathtools}
\usepackage{bm}
% indicator function via boldsymbol (no extra package needed)
\usepackage{booktabs}
\usepackage{graphicx}
\usepackage{hyperref}
\usepackage[margin=2.5cm]{geometry}
\usepackage{enumitem}
\usepackage{xcolor}
\usepackage{tcolorbox}
\usepackage{float}

% Theorem environments
\theoremstyle{definition}
\newtheorem{definition}{Definition}[section]
\newtheorem{example}{Example}[section]
\theoremstyle{plain}
\newtheorem{theorem}{Theorem}[section]
\newtheorem{proposition}{Proposition}[section]
\theoremstyle{remark}
\newtheorem{remark}{Remark}[section]

% Colored boxes
\newtcolorbox{keyresult}[1][]{
    colback=blue!5!white,
    colframe=blue!60!black,
    fonttitle=\bfseries,
    title={Key Insight},
    #1
}
\newtcolorbox{warningbox}[1][]{
    colback=red!5!white,
    colframe=red!60!black,
    fonttitle=\bfseries,
    title={Warning},
    #1
}
\newtcolorbox{refbox}[1][]{
    colback=gray!5!white,
    colframe=gray!50!black,
    fonttitle=\bfseries,
    title={Go Deeper},
    #1
}

% Custom commands
\newcommand{\VaR}{\operatorname{VaR}}
\newcommand{\ES}{\operatorname{ES}}
\newcommand{\CVaR}{\operatorname{CVaR}}
\newcommand{\EV}{\mathbb{E}}
\newcommand{\PV}{\mathbb{P}}
\newcommand{\RV}{\mathbb{R}}
\newcommand{\ind}{\boldsymbol{1}}

\title{%
    \textbf{Foundations of Modern Financial Risk} \\[0.5em]
    \large What You Need to Know Before Writing Code
}
\author{Risk Analyst Project}
\date{March 2026}

\begin{document}
\maketitle

\begin{abstract}
A concise, example-driven introduction to the core concepts of modern financial risk analysis.
Each idea is motivated by a concrete scenario, then generalized.
We assume familiarity with probability, statistics, and basic stochastic processes;
the focus is on \emph{how these tools are used in risk management} and \emph{why the field evolved the way it did}.
References point to deeper material for each topic.
\end{abstract}

\tableofcontents
\newpage

%% ============================================================================
\section{What Risk Means in Finance}
%% ============================================================================

Risk in finance is not synonymous with volatility.
It is the \emph{possibility of loss} --- specifically, loss that exceeds what you planned for.

\begin{example}[A simple portfolio]
You hold \$100{,}000 split equally between the S\&P~500 (SPY) and 20-year Treasury bonds (TLT).
On a typical day the portfolio moves $\pm$\$500.
But on March~16, 2020, the S\&P dropped 12\% in a single session while Treasuries also sold off.
Your portfolio lost \$8{,}200 --- sixteen times the ``normal'' day.

Risk management answers: \emph{How often can losses this large happen?
How much capital should I reserve?
What if it gets even worse?}
\end{example}

The field has converged on three interlocking concerns:

\begin{enumerate}[leftmargin=2cm]
    \item[\textbf{Measure}] Quantify the loss distribution: how bad can it get, and with what probability?
    \item[\textbf{Manage}] Set limits, allocate capital, hedge, diversify.
    \item[\textbf{Monitor}] Backtest your models against reality; adjust when they fail.
\end{enumerate}

%% ============================================================================
\section{Loss Distributions and Tail Behavior}
%% ============================================================================

\subsection{From Returns to Losses}

Let $r_t = \ln(P_t / P_{t-1})$ be the log-return of an asset.
For a portfolio with weights $\bm{w}$ and asset returns $\bm{r}_t$, the portfolio return is $r_p = \bm{w}^\top \bm{r}_t$.
The \textbf{loss} is $L = -r_p$ (we flip the sign so that losses are positive).

The central object in risk is the \textbf{loss distribution} $F_L(x) = \PV(L \leq x)$.

\subsection{Why the Normal Distribution Fails}

\begin{example}[Fat tails in practice]\label{ex:fat-tails}
Daily S\&P~500 returns from 1990--2025 have:
\begin{itemize}
    \item Kurtosis $\approx 11.5$ (normal = 3). Tails are \emph{much} heavier.
    \item The Oct~19, 1987 crash ($-22.6\%$) is a $\sim$25-sigma event under normality --- probability $\approx 10^{-138}$.
    \item In reality, moves $> 4\sigma$ happen about 10x more often than the normal predicts.
\end{itemize}
\end{example}

\begin{keyresult}
Financial returns exhibit \textbf{heavy tails} (excess kurtosis), \textbf{volatility clustering} (large moves follow large moves), and \textbf{asymmetry} (crashes are sharper than rallies).
Any risk model that ignores these features will systematically underestimate losses.
\end{keyresult}

Better distributional choices:
\begin{itemize}
    \item \textbf{Student-$t$}: $f(x; \nu) \propto (1 + x^2/\nu)^{-(\nu+1)/2}$. Tails decay polynomially, not exponentially. Typical fit: $\nu \approx 4$--$6$ for equity returns.
    \item \textbf{Generalized Pareto} (for extreme tails only): the correct asymptotic model for tail exceedances (see Section~\ref{sec:evt}).
    \item \textbf{Skewed-$t$} (Hansen, 1994): adds asymmetry to Student-$t$. Used as the innovation distribution in GARCH models.
\end{itemize}

\begin{refbox}
\begin{itemize}
    \item Cont, R. (2001). ``Empirical properties of asset returns: stylized facts and statistical issues.'' \emph{Quantitative Finance}, 1(2), 223--236. \textit{The} survey of return properties.
    \item McNeil, Frey \& Embrechts (2015). \emph{Quantitative Risk Management}, Ch.~3. Distributional models for returns.
\end{itemize}
\end{refbox}

%% ============================================================================
\section{Risk Measures: VaR, Expected Shortfall, and Coherence}
\label{sec:risk-measures}
%% ============================================================================

\subsection{Value-at-Risk}

\begin{definition}[Value-at-Risk]
For a loss $L$ with distribution $F_L$ and confidence level $\alpha \in (0,1)$:
\[
    \VaR_\alpha(L) = F_L^{-1}(\alpha) = \inf\{x : F_L(x) \geq \alpha\}
\]
\end{definition}

\begin{example}[Computing VaR]
Your SPY/TLT portfolio has 252 daily returns over the past year.
Sort the losses from smallest to largest.
At $\alpha = 0.95$ (95\%), $\VaR_{0.95}$ is the 13th-worst loss (since $252 \times 0.05 \approx 13$).

Suppose that value is \$1{,}850.
Interpretation: ``On 95\% of days, losses do not exceed \$1{,}850.''
Equivalently: ``We expect to lose more than \$1{,}850 about once every 20 trading days.''
\end{example}

\textbf{Three ways to compute VaR:}
\begin{enumerate}
    \item \textbf{Historical simulation}: sort historical losses, read off the quantile.
    \item \textbf{Parametric (variance-covariance)}: assume $L \sim N(\mu, \sigma^2)$, then $\VaR_\alpha = \mu + \sigma \, \Phi^{-1}(\alpha)$. Fast but relies on normality.
    \item \textbf{Monte Carlo}: simulate thousands of portfolio paths, compute losses, read off the quantile. Most flexible; handles non-linear instruments.
\end{enumerate}

\subsection{The Problem with VaR}

\begin{warningbox}
VaR tells you the \emph{threshold} but nothing about \emph{how bad it gets beyond that threshold}.
Two portfolios can have identical VaR but very different tail losses.
\end{warningbox}

\begin{example}[VaR blindness]
Portfolio A loses exactly \$2{,}000 on the 5\% worst days.
Portfolio B loses \$2{,}000 on most of those days, but occasionally loses \$50{,}000.
Both have $\VaR_{0.95} = \$2{,}000$.
A risk manager using only VaR would treat them identically.
\end{example}

Worse: VaR is \textbf{not subadditive} in general.
It is possible that $\VaR(A+B) > \VaR(A) + \VaR(B)$, meaning VaR can penalize diversification.

\subsection{Expected Shortfall (CVaR)}

\begin{definition}[Expected Shortfall]
\[
    \ES_\alpha(L) = \frac{1}{1-\alpha} \int_\alpha^1 \VaR_u(L) \, du
\]
When $F_L$ is continuous, this simplifies to:
\[
    \ES_\alpha(L) = \EV[L \mid L > \VaR_\alpha(L)]
\]
\end{definition}

ES is the \emph{average loss in the tail} beyond VaR.
It directly answers: ``When things go wrong, how wrong do they go on average?''

\begin{example}[ES vs.\ VaR]
For the same SPY/TLT portfolio, if $\VaR_{0.95} = \$1{,}850$, then $\ES_{0.95}$ averages the 13 worst losses.
Suppose those are \$1{,}850, \$2{,}100, \$2{,}300, \ldots, \$8{,}200.
Their average might be \$3{,}400.
So $\ES_{0.95} = \$3{,}400$ --- nearly double the VaR.
\end{example}

\subsection{Coherent Risk Measures}

Artzner et al.\ (1999) defined four axioms a ``sensible'' risk measure $\rho$ should satisfy:

\begin{enumerate}
    \item \textbf{Monotonicity}: If $L_1 \leq L_2$ always, then $\rho(L_1) \leq \rho(L_2)$.
    \item \textbf{Subadditivity}: $\rho(L_1 + L_2) \leq \rho(L_1) + \rho(L_2)$. Diversification does not increase risk.
    \item \textbf{Positive homogeneity}: $\rho(\lambda L) = \lambda \, \rho(L)$ for $\lambda > 0$.
    \item \textbf{Translation invariance}: $\rho(L + c) = \rho(L) + c$. Adding cash reduces risk dollar-for-dollar.
\end{enumerate}

\begin{keyresult}
\textbf{ES is coherent. VaR is not} (it fails subadditivity).
This is the theoretical reason Basel III/IV replaced VaR with ES for market risk capital:
the Fundamental Review of the Trading Book (FRTB) mandates $\ES_{0.975}$ instead of $\VaR_{0.99}$.
\end{keyresult}

\begin{refbox}
\begin{itemize}
    \item Artzner, Delbaen, Eber \& Heath (1999). ``Coherent measures of risk.'' \emph{Mathematical Finance}, 9(3), 203--228. The foundational paper.
    \item Acerbi \& Tasche (2002). ``On the coherence of expected shortfall.'' \emph{Journal of Banking \& Finance}, 26(7), 1487--1503.
    \item Jorion (2007). \emph{Value at Risk}. The practitioner's reference.
\end{itemize}
\end{refbox}

%% ============================================================================
\section{Extreme Value Theory: Modeling the Worst Cases}
\label{sec:evt}
%% ============================================================================

Standard distributional fits (normal, $t$) apply to the \emph{bulk} of the distribution.
For the extreme tails --- the 1\% or 0.1\% worst outcomes --- Extreme Value Theory (EVT) provides a rigorous framework.

\subsection{The Key Theorem}

\begin{theorem}[Pickands--Balkema--de Haan]
For a wide class of distributions, the excess loss above a high threshold $u$ converges in distribution to the \textbf{Generalized Pareto Distribution} (GPD):
\[
    \PV(L - u \leq y \mid L > u) \approx G_{\xi, \beta}(y) =
    \begin{cases}
        1 - (1 + \xi y / \beta)^{-1/\xi} & \xi \neq 0 \\
        1 - e^{-y/\beta} & \xi = 0
    \end{cases}
\]
where $\xi$ is the \textbf{shape} (tail index) and $\beta > 0$ is the scale.
\end{theorem}

\begin{example}[Tail index interpretation]
\begin{itemize}
    \item $\xi > 0$ (heavy tail): Pareto-like. Typical for equity returns ($\xi \approx 0.2$--$0.4$).
    \item $\xi = 0$ (light tail): exponential decay. Rare in finance.
    \item $\xi < 0$ (bounded tail): finite upper endpoint. Insurance losses sometimes.
\end{itemize}
Fitting GPD to S\&P daily losses above the 95th percentile typically gives $\hat\xi \approx 0.25$, confirming heavy tails and allowing extrapolation to extreme quantiles (99.9\%) that have few or no historical observations.
\end{example}

\begin{refbox}
\begin{itemize}
    \item de Haan \& Ferreira (2006). \emph{Extreme Value Theory: An Introduction}. The mathematical reference.
    \item McNeil \& Frey (2000). ``Estimation of tail-related risk measures for heteroscedastic financial time series.'' \emph{Journal of Empirical Finance}. The GARCH + EVT approach.
    \item Embrechts, Kl\"uppelberg \& Mikosch (1997). \emph{Modelling Extremal Events}. Classic for insurance/finance.
\end{itemize}
\end{refbox}

%% ============================================================================
\section{Dependence: Beyond Correlation}
%% ============================================================================

\subsection{Why Correlation Is Not Enough}

Pearson correlation $\rho$ measures \emph{linear} association.
Two variables can have $\rho = 0$ yet be strongly dependent (e.g., $Y = X^2$).
More critically for risk: assets that appear weakly correlated in normal markets can become highly correlated during crises.

\begin{example}[Correlation breakdown]
From 2005--2007, the correlation between US equities and mortgage-backed securities was $\approx 0.3$.
During the 2008 crisis, it spiked above $0.8$.
A portfolio diversified under normal-market correlation was far less diversified when it mattered most.
\end{example}

\subsection{Copulas: The Right Tool}

A \textbf{copula} $C(u_1, \ldots, u_d)$ is a function that binds marginal distributions into a joint distribution.
By Sklar's theorem, \emph{any} joint distribution can be decomposed as:
\[
    F(x_1, \ldots, x_d) = C\bigl(F_1(x_1), \ldots, F_d(x_d)\bigr)
\]
This separates two modeling choices:
\begin{itemize}
    \item \textbf{Marginals} $F_i$: model each asset's return distribution individually (e.g., GARCH with skewed-$t$).
    \item \textbf{Dependence structure} $C$: model how they move together, especially in the tails.
\end{itemize}

Key copula families and their tail behavior:

\begin{table}[H]
\centering
\begin{tabular}{@{}lcc@{}}
\toprule
\textbf{Copula} & \textbf{Lower tail dep.}\ $\lambda_L$ & \textbf{Upper tail dep.}\ $\lambda_U$ \\
\midrule
Gaussian       & 0 & 0 \\
Student-$t$    & $> 0$ (symmetric) & $> 0$ (symmetric) \\
Clayton        & $> 0$ & 0 \\
Gumbel         & 0 & $> 0$ \\
Frank          & 0 & 0 \\
\bottomrule
\end{tabular}
\caption{Tail dependence by copula family. For risk, lower tail dependence (joint crashes) matters most. The Gaussian copula has \emph{none} --- this was a core flaw in pre-2008 CDO pricing.}
\end{table}

\begin{keyresult}
The Gaussian copula assigns probability zero to joint extreme co-movements.
This made it catastrophically wrong for pricing CDOs and estimating portfolio credit risk before 2008.
Modern practice uses $t$-copulas or Clayton copulas that capture tail dependence.
\end{keyresult}

\begin{refbox}
\begin{itemize}
    \item Nelsen (2006). \emph{An Introduction to Copulas}. 2nd ed. The textbook.
    \item Joe (2014). \emph{Dependence Modeling with Copulas}. Advanced, covers vine copulas.
    \item Embrechts, McNeil \& Straumann (2002). ``Correlation and dependence in risk management.'' \emph{Risk Magazine}. Why correlation fails.
\end{itemize}
\end{refbox}

%% ============================================================================
\section{The Regulatory Landscape: Why This Matters}
%% ============================================================================

Risk models do not exist in a vacuum.
Banks, insurers, and asset managers operate under regulatory frameworks that dictate \emph{which} risk measures to compute, \emph{how} to compute them, and \emph{how much capital} to hold.

\subsection{Basel Framework (Banking)}

\begin{table}[H]
\centering
\begin{tabular}{@{}llp{7.5cm}@{}}
\toprule
\textbf{Accord} & \textbf{Era} & \textbf{Key change for risk measurement} \\
\midrule
Basel I (1988)   & Simple & Flat capital charges by asset class. No internal models. \\
Basel II (2004)  & VaR era & Banks allowed to use internal VaR models (99\%, 10-day). \\
Basel III (2010) & Post-crisis & Stressed VaR, capital buffers, leverage ratio added. \\
Basel III.1 / ``IV'' (2017--2028) & Current & \textbf{FRTB}: replaces VaR with ES at 97.5\%. Output floors on internal models. Standardized Measurement Approach for operational risk. \\
\bottomrule
\end{tabular}
\end{table}

\textbf{FRTB (Fundamental Review of the Trading Book)}:
\begin{itemize}
    \item Market risk capital: $\ES_{0.975}$ over a base horizon of 10 days, computed with different liquidity horizons per risk factor.
    \item Internal Models Approach (IMA): banks must pass backtesting at the desk level.
    \item Standardized Approach (SA): sensitivities-based method (delta, vega, curvature) for banks without IMA approval.
    \item Implementation: EU phasing 2025--2026, UK targeted 2027, US ``Basel Endgame'' from July 2025.
\end{itemize}

\subsection{What This Means for You}

Every project in this repo connects to real regulatory or industry requirements:
\begin{itemize}
    \item \textbf{P01} (Dashboard): implements VaR and ES with Basel-mandated backtesting.
    \item \textbf{P02} (Monte Carlo): the computational engine behind IMA.
    \item \textbf{P03} (Credit Scoring): PD models under Basel IRB, with SR~11-7 model governance.
    \item \textbf{P04} (Volatility): GARCH models feed into ES and VaR calculations.
\end{itemize}

\begin{refbox}
\begin{itemize}
    \item Basel Committee (2019). ``Minimum capital requirements for market risk.'' BCBS d457. The FRTB standard.
    \item Federal Reserve (2011). SR~11-7: ``Guidance on Model Risk Management.'' Gold standard for model governance.
    \item Hull (2018). \emph{Risk Management and Financial Institutions}. Ch.~15--17 on Basel.
\end{itemize}
\end{refbox}

%% ============================================================================
\section{Stochastic Processes for Risk}
%% ============================================================================

You already know stochastic processes from your actuarial training.
Here we focus on the specific models used in financial risk and \emph{why} each matters.

\subsection{Geometric Brownian Motion (GBM)}

The workhorse of Monte Carlo simulation:
\[
    dS_t = \mu S_t \, dt + \sigma S_t \, dW_t
    \quad \Longrightarrow \quad
    S_T = S_0 \exp\!\left[\left(\mu - \tfrac{\sigma^2}{2}\right)T + \sigma W_T\right]
\]
\textbf{Used in}: option pricing (Black-Scholes), portfolio simulation, VaR by Monte Carlo.

\textbf{Limitation}: constant volatility $\sigma$.
Markets exhibit volatility clustering --- addressed by GARCH (discrete) and stochastic volatility models (continuous).

\subsection{GARCH(1,1): Volatility That Remembers}

\[
    r_t = \mu + \epsilon_t, \quad \epsilon_t = \sigma_t z_t, \quad z_t \sim D(0,1)
\]
\[
    \sigma_t^2 = \omega + \alpha \epsilon_{t-1}^2 + \beta \sigma_{t-1}^2
\]
where $D$ is typically Student-$t$ or skewed-$t$.

\begin{example}[GARCH intuition]
If yesterday's return was large ($|\epsilon_{t-1}|$ big), today's conditional variance $\sigma_t^2$ increases.
The $\beta$ term gives persistence: high volatility regimes last for weeks or months.
For the S\&P~500, typical estimates are $\alpha \approx 0.09$, $\beta \approx 0.90$, giving persistence $\alpha + \beta \approx 0.99$ --- volatility shocks decay very slowly.
\end{example}

\textbf{Extensions} (covered in Project~04): GJR-GARCH (leverage effect: negative returns increase volatility more), EGARCH (log-volatility, no positivity constraint), APARCH (power transformations).

\subsection{Mean-Reverting Processes}

For interest rates and credit spreads, the Ornstein-Uhlenbeck / Vasicek model:
\[
    dr_t = \kappa(\theta - r_t)\,dt + \sigma\,dW_t
\]
Mean-reverts to $\theta$ at speed $\kappa$.
Used in: yield curve modeling, CVA calculations (Project~09), credit intensity models.

\begin{refbox}
\begin{itemize}
    \item Glasserman (2003). \emph{Monte Carlo Methods in Financial Engineering}. Ch.~3 (GBM simulation), Ch.~8 (variance reduction).
    \item Engle (1982). ``Autoregressive conditional heteroscedasticity.'' \emph{Econometrica}. The ARCH origin.
    \item Bollerslev (1986). ``Generalized autoregressive conditional heteroscedasticity.'' \emph{Journal of Econometrics}. GARCH(1,1).
\end{itemize}
\end{refbox}

%% ============================================================================
\section{Summary: The Mental Model}
%% ============================================================================

\begin{center}
\begin{tabular}{@{}p{3cm}p{4.5cm}p{5cm}@{}}
\toprule
\textbf{Layer} & \textbf{What it does} & \textbf{Key tools} \\
\midrule
Marginals & Model each risk factor's distribution & GARCH, Student-$t$, skewed-$t$, EVT (GPD) \\
\addlinespace
Dependence & Model how risk factors move together & Copulas (Gaussian, $t$, Clayton, vine) \\
\addlinespace
Aggregation & Combine into portfolio loss distribution & Monte Carlo, variance reduction, FFT \\
\addlinespace
Risk measures & Summarize the loss distribution & VaR, ES, spectral measures \\
\addlinespace
Validation & Check if the model works & Backtesting (Kupiec, Christoffersen), stress testing \\
\bottomrule
\end{tabular}
\end{center}

Every project in the curriculum operates within this pipeline.
The first four projects build it from the ground up:
the dashboard (P01) computes measures and validates them;
the MC engine (P02) handles aggregation;
credit scoring (P03) models a different kind of marginal (default probability);
and volatility modeling (P04) improves the marginal layer.

\vfill
\hrule
\smallskip
\noindent\textit{Risk Analyst Project} \hfill \textit{March 2026}

\end{document}
